\documentclass[12pt]{article}
\usepackage[utf8]{inputenc}
\usepackage[T1]{fontenc}
\usepackage[margin=1in]{geometry}
\usepackage{setspace}
\usepackage{xcolor}
\usepackage[colorlinks=true, urlcolor=blue!60!black, linkcolor=blue!60!black]{hyperref}
\setlength{\parindent}{0em}
\setlength{\parskip}{0.6em}
\renewcommand{\baselinestretch}{1.0}
\begin{document}
\begin{center}
{\LARGE \textit{Poindexter v. Greenhow}, 114 U.S. 270 (1885)\par}
\vspace{0.5em}
{\large Case Brief --- Two-Page Edition\par}
\end{center}

\section{Memory Jogger}
In the lead \textit{Virginia Coupon Case}, a 5--4 Court held that Virginia's tax-receivable bond coupons were valid contracts, not bills of credit; that refusing them impaired the obligation of contracts; and that a taxpayer could sue the collector who seized his property, because an officer enforcing a void statute is not the State (\href{https://supreme.justia.com/cases/federal/us/114/270/}{114 U.S. 270}).

\section{Detailed Case Facts}
Virginia's 1871 Funding Act issued bonds whose coupons would ``be receivable at and after maturity for all taxes, debts, dues, and demands due the state.'' An 1882 Act forbade collectors to accept them and abolished common-law remedies against collectors. Richmond treasurer Samuel Greenhow refused Poindexter's coupons for \$12.45 in taxes and seized his office desk.

\section{Procedural History}
Poindexter sued Greenhow personally in detinue. The Hustings Court of Richmond ruled for Greenhow, and the Supreme Court took the case with seven companion coupon cases.

\section{Judicial Votes}
Reversed, 5--4. Matthews wrote for Field, Harlan, Woods, and Blatchford; Bradley, with Waite, Miller, and Gray, dissented.

\section{Holding}
The coupons are not bills of credit; the 1882 Act unconstitutionally impaired Virginia's contract and was wholly void; and because tender equals payment, Greenhow acted without authority and could be sued personally despite the Eleventh Amendment.

\section{Analysis of Opinions}
Matthews reasoned that bills of credit are designed to circulate as money, while the coupons were receivable only at maturity. Tax receivability was the essence of the contract, and the pay-under-protest remedy destroyed it. Distinguishing suits to compel payment from the treasury, barred by \textit{Louisiana v. Jumel}, from suits against an officer who wrongs a citizen's property, he held that a void statute is ``as if it were not and never had been,'' leaving Greenhow a trespasser. Bradley called this a fiction: any suit designed to coerce a State's performance is a suit against the State. \textit{Ex parte Young} later relied on \textit{Poindexter}.

\section{Examples: Future Applications}
\textbf{Same-Side:} a company may sue the revenue officer who refused its coupons and seized its truck. \textbf{Opposite-Side:} a suit ordering the treasurer to pay overdue coupons from the treasury is barred under \textit{Jumel} and \textit{In re Ayers}. \textbf{Fence-Sitter:} an injunction requiring collectors to accept coupons in the future.

\section{Critique}
The majority treats Greenhow as not acting for Virginia for immunity purposes, yet his conduct violates the Contracts Clause only because he enforces Virginia's law. That paradox passed into \textit{Ex parte Young} as the defining fiction of federal courts doctrine.

\section{Key Quotations}
\begin{quote}
``The defendant in error is not her officer, her agent, or her representative, in the matter complained of; for he has acted not only without her authority, but contrary to her express commands.''
\end{quote}
\noindent --- Justice Matthews, Opinion of the Court, 114 U.S. at 292--93.

\end{document}
